\subsection{Codebase Organization}
\label{sec:code-organization}

The source tree separates operating-system integration, reusable numeric and
rendering code, application front ends, and validation tools.  The build-target
boundaries are significant: they keep CUDA compilation isolated, prevent
platform APIs from leaking into portable interfaces, and allow the Windows GUI,
Linux GUI, and command-line renderer to share the same rendering implementation.

\subsubsection{Dependency layers}

\Cref{fig:codebase-layers} summarizes the architectural layers.  Arrows point
toward targets that consume the layer below them; final executables link the
libraries required by the features they expose.

\begin{figure}[!htbp]
  \centering
  \begin{tikzpicture}[
      component/.style={
        draw,
        rounded corners,
        align=center,
        minimum height=12mm,
        text width=0.90\textwidth,
        inner sep=3mm
      },
      flow/.style={-{Latex[length=2.5mm]}, thick}
    ]
    \node[component] (foundation) {
      \textbf{Foundation libraries}\\
      \code{FractalSharkPlatform}\qquad \code{HpSharkFloatLib}
    };
    \node[component, above=8mm of foundation] (rendering) {
      \textbf{Shared rendering implementation}\\
      \code{FractalSharkLib}\qquad \code{FractalSharkGpuLib}
    };
    \node[component, above=8mm of rendering] (consumers) {
      \textbf{Applications and developer tools}\\[1mm]
      \code{FractalShark} + \code{FractalSharkGuiWin32}\quad
      \code{FractalSharkGuiLinux}\quad \code{FractalSharkCli}\\
      \code{FractalTray}\quad \code{FractalSaver}\quad
      \code{FractalSharkTest}\quad \code{HpSharkFloatTest}\\
      \code{HpSharkFloatTestLib}\quad \code{HpSharkInstantiate}
    };

    \draw[flow] (foundation) -- (rendering);
    \draw[flow] (rendering) -- (consumers);
  \end{tikzpicture}
  \caption[Codebase dependency layers]{Codebase dependency layers.  Platform and
    high-precision arithmetic support the shared CPU and CUDA rendering code;
    applications and validation tools consume those shared layers.}
  \label{fig:codebase-layers}
\end{figure}

\subsubsection{Foundation and rendering libraries}

\paragraph{\code{FractalSharkPlatform}.}
The platform library owns operating-system services used by otherwise portable
code.  Its common interfaces cover allocation, files, process information,
timing, diagnostics, crash handling, wait-cursor state, process memory limits,
and native OpenGL contexts.  Separate Win32 and Linux translation units provide
the implementations described in \cref{subsec:platform-abstraction}.

\paragraph{\code{HpSharkFloatLib}.}
The high-precision GPU arithmetic library implements the limb-based
\code{HpSharkFloat} representation, addition, NTT multiplication,
normalization, periodicity support, and the GPU reference-orbit machinery.  It
also contains the parameter families used to specialize those kernels.  The
numeric representation and execution pipeline are documented in
\cref{subsec:ref-orbit-gpu,subsec:hpfloat-model,sec:hp-add,sec:ntt-multiply}.

\paragraph{\code{FractalSharkLib}.}
The core library owns portable fractal state and CPU-side orchestration.  Its
responsibilities include coordinate conversion, reference-orbit dispatch,
linear-approximation construction, feature finding, palettes, iteration-buffer
management, render-thread coordination, persistence, PNG output, shared
commands, and functional tests.  The principal flows appear in
\cref{sec:ref-orbit-calc,sec:render-thread-pool,sec:disk_backed_growable_vectors}.

\paragraph{\code{FractalSharkGpuLib}.}
The rendering CUDA library owns the per-pixel launch surface and kernels for
direct escape-time rendering, perturbation, BLA and LA~v2 traversal,
antialiasing, coloring, and reductions.  It consumes numeric and state types
defined by the core and high-precision libraries; final applications link all
required libraries together.  Kernel behavior is covered in
\cref{sec:mandel-base-kernels,sec:perturb-only,sec:la-v2-perturb,sec:kernel-list}.

\subsubsection{Applications and front ends}

\paragraph{\code{FractalShark} and \code{FractalSharkGuiWin32}.}
The primary Windows application is split across two Visual~Studio projects.
\code{FractalShark} supplies the executable target, resources, icons, and final
link, while \code{FractalSharkGuiWin32} implements the Win32 entry point, main
window, command dispatcher, popup menus, console, saved-location UI, and splash
window.  Shared command identifiers remain in the core library.  User-visible
behavior and frame presentation are described in
\cref{sec:ui-popup-menu,sec:keyboard-shortcuts,subsec:opengl-display}.

\paragraph{\code{FractalSharkGuiLinux}.}
The native Linux application combines an Xlib event loop and menus, GLX context
management, Dear~ImGui overlays, clipboard integration, and the shared renderer.
Its platform boundary and current parity limitations are documented in
\cref{subsec:linux-gui}.

\paragraph{\code{FractalSharkCli}.}
The portable command-line executable provides headless rendering from built-in
views, saved locations, or explicit coordinates.  It exercises the same core
and CUDA rendering libraries without a window system; its contract and render
path are documented in \cref{sec:fractalshark-cli}.

\paragraph{\code{FractalTray}.}
The Windows tray application batch-renders a sequence of coordinates through
the shared renderer.  Its input format, configuration, and operating procedure
are described in \cref{subsec:fractaltray}.

\paragraph{\code{FractalSaver}.}
The Windows screen-saver project is legacy, unsupported code retained in the
Visual~Studio solution.  It is not part of the CMake build or release packages
and should not be treated as an active front end.

\subsubsection{Validation and generation tools}

\paragraph{\code{FractalSharkTest}.}
The portable unit-test executable covers core numeric types, coordinate
conversion, linear-approximation data, allocators, serialization, command
state, render goldens, and related utilities.  The custom test framework and
representative coverage are summarized in \cref{subsec:fractalsharktest}.

\paragraph{\code{HpSharkFloatTestLib} and \code{HpSharkFloatTest}.}
The support library contains reference implementations, correctness tests,
performance drivers, and result tracking for high-precision CUDA arithmetic.
The executable supplies the CUDA-aware harness that invokes those tests.
Runtime validation requires physical CUDA hardware, as discussed in
\cref{subsec:hpsharkfloattest}.

\paragraph{\code{HpSharkInstantiate}.}
The standalone developer tool generates explicit-instantiation CUDA sources
for the fixed \code{SharkFloatParams} families.  Generated translation units
control template compile time and final binary coverage; the underlying
strategy is documented in \cref{subsec:template-explicit-inst}.

\subsubsection{Build-system coverage}

The Visual~Studio solution contains the Windows application, Windows GUI
library, tray application, legacy screen saver, shared libraries, tests, and
generation tool.  CMake builds the portable shared libraries and tools,
\code{FractalSharkCli}, and \code{FractalSharkGuiLinux}; it intentionally omits
the Win32-only front ends.  Parallel project-file maintenance, Linux target
composition, and continuous-integration coverage are described in
\cref{subsec:dual-build}.
